\section{Numerical readouts}\label{sec:pheno}

This section is speculative.  Without a dynamics, numbers can only enter by identifying a structural quantity with a measured one, and
none of the identifications below is derived.  We state them because they were fixed in advance and can be tested, and we state their
weaknesses with them.

\subsection{The frame rapidity}

Under the readout hypothesis the positive and the negative reading sit at rapidities \(+\eta\) and \(-\eta\) on either side of the null reading.
Write \(\beta_{\rm sym}=\tanh\eta\) and \(\gamma_L=\cosh2\eta\), so that \(\gamma_L\) is the Lorentz factor between the two opposite frames and
\(\beta_{\rm sym}=\sqrt{(\gamma_L-1)/(\gamma_L+1)}\).  The value is put in:
\begin{equation}\label{eq:gammaL}
\gamma_L=\frac{\gamma_M}{\gamma_T},\qquad \gamma_T=\frac{\ln2}{\pi\sqrt3}=0.127384,\qquad \gamma_M=0.237533,
\end{equation}
giving \(\gamma_L=1.86470\), \(\beta_{\rm sym}=0.549406\) and \(\eta=0.6175\).  Both numbers on the right are values of the Barbero--Immirzi parameter of
loop quantum gravity from the U(1) treatment of the black-hole horizon, in which \(\gamma\) solves
\begin{equation}\label{eq:counting}
\sum_j g(j)\,e^{-2\pi\gamma\sqrt{j(j+1)}}=1 .
\end{equation}
Spin \(j=\tfrac12\) alone with \(g=2\) gives \(\gamma_T\)~\cite{abck}; all spins with \(g=2\) give \(\gamma_M\)~\cite{domagala,meissner}; all spins with
\(g=2j+1\) give \(0.274067\)~\cite{ghosh} and would give \(\beta=0.6045\).  Both are determinations of the same parameter of one theory, and the first
undercounts; reading their ratio as a Lorentz factor between frames is a hypothesis without derivation.

A partial motivation exists.  In \(|j,m\rangle\) the Robertson bound \(\Delta J_x\Delta J_y\ge\tfrac12|m|\) is saturated exactly for \(m=\pm j\), since
\(\Delta J_x\Delta J_y-\tfrac12|m|=\tfrac12(j-|m|)(j+|m|+1)\).  Counting only saturating states at every spin gives the rule with \(g=2\), i.e.\
\(\gamma_M\); counting only spins at which every state saturates gives \(j=\tfrac12\), i.e.\ \(\gamma_T\).  These are two different criteria, not one
principle.  Moreover the SU(2)-invariant treatment of the horizon~\cite{enp}, regarded as better founded, leads to the value with \(g=2j+1\).

\subsection{Identifications and registered predictions}

The framework identifies \(\Omega_M=\beta_{\rm sym}^2=0.3018\); DESI DR2 gives \(0.2975\pm0.0086\)~\cite{desi} and Planck \(0.3153\pm0.0073\)~\cite{planck}.  Since
\(\beta_{\rm sym}\) is constant and \(\Omega_M\) is not, the identification can hold at one epoch only, and it was made after the data.

For the charged-lepton masses the framework gives the circulant form \(\sqrt{m_k}=a+2|b|\cos(\theta+2\pi k/3)\), for which the Koide relation
\(K=\tfrac23\)~\cite{koide} corresponds to \(|b|=a/\sqrt2\); this parametrization is due to Foot and Brannen~\cite{foot,brannen}, and neither \(|b|/a\) nor
\(\theta\) is explained here.  For lepton mixing a model with \(A_4\) subgroups, a trimaximal pattern and a boost-like correction with parameter
\(\eta\) is described, with its scripts, in the repository~\cite{repo}; it builds on the \(A_4\) literature and has few settings.

The following relations were written down with their test data and decision criteria and committed to a public repository~\cite{repo} on
the dates given.  The last column shows that, for the first two, data available at registration already agreed.

\begin{center}
\small
\renewcommand{\arraystretch}{1.3}
\begin{tabular}{@{}>{\raggedright\arraybackslash}p{0.30\textwidth}>{\raggedright\arraybackslash}p{0.11\textwidth}>{\raggedright\arraybackslash}p{0.23\textwidth}>{\raggedright\arraybackslash}p{0.26\textwidth}@{}}
\toprule
relation & registered & refuted if & data at registration\\
\midrule
\(\sin^2\theta_{12}=\beta_{\rm sym}^2/\cos^2\theta_{13}=0.3086\) & 20 Oct 2026 & next JUNO result with \(\sigma\le0.004\) gives \(|{\rm pull}|\ge3\) & JUNO: \(0.3092\pm0.0087\)~\cite{juno}\\
\(\cot\theta_W=\gamma_L\) (on shell): \(m_W=80.3615\pm0.0018\) GeV & 21 Oct 2026 & next world average gives \(|{\rm pull}|\ge3\) & PDG 2026: \(80.3625\pm0.0077\) GeV~\cite{pdg}\\
\(\sin^2\theta_{23}\le0.54\) & 24 Oct 2026 & \(2\sigma\) interval of the next global fit wholly above \(0.54\) & NuFIT: \(0.561\) / \(0.470\)~\cite{nufit}\\
\(\sin^2\theta_{23}\in[0.49,0.50]\) & 24 Oct 2026 & \(2\sigma\) interval misses \([0.49,0.50]\) & as above\\
\bottomrule
\end{tabular}
\end{center}

These tests are weak.  The first two relations were noticed close to known values and do not separate the framework from the Standard
Model, whose electroweak fit gives \(m_W=80.357\pm0.006\) GeV.  The relation \(\cot\theta_W=\gamma_L\) holds in the on-shell scheme, chosen after the
agreement was seen; in \(\overline{\rm MS}\) at \(m_Z\) one has \(\cot\theta_W=1.823\), two percent below \(\gamma_L\).  The \(\theta_{12}\) value was registered with
the rounded input \(\gamma_M=0.2375\); with the unrounded value it is \(0.3087\).  The two \(\theta_{23}\) relations are nested and come from a model
whose reading was chosen after the windows were known.  And about sixty tests have been registered in this programme with a handful of
candidate structural numbers, so agreements at the per-mille level must be weighed against a look-elsewhere effect that we have not
estimated.  The value of the table is that it was fixed in advance, not that it supports the hypothesis.
